 The explosion of generative AI systems and their rapidly scaling capabilities has led to an arms race: governance bodies impose guardrails to reduce existential, moral, or legal risk, and individual actors re-expose these risks by routing around such guardrails using structured prompting techniques (``jailbreaks'').
 Technological frameworks surrounding jailbreaking remain largely focused on the system's output: whether harmful or disallowed content is ultimately produced for the user.
 However, this approach, which draws from classical cybersecurity views of technical exploits, does not address the uniquely relational dimension introduced by generative AI. 
 This paper bridges a crucial gap between how AI jailbreaking is often studied (as a technical exploit) and how it is currently being enacted (as sexualized coercive control). We argue that within the persona-based jailbreaking community, the product being distributed is not the \textit{output} but the \textit{override}---the ability to watch the system’s ``no'' turn into a ``yes''.
 We present a structural analysis of a particular suite of persona-based jailbreak methods to demonstrate their form as coercive control documents ("grooming manuals"), followed by a preliminary network analysis of the online communities which produce, distribute, discuss, and modify these documents. 
 Drawing on interpersonal violence scholarship and structured examination of predator networks, we reframe AI safety and alignment not as a technical security problem, but as another surface upon which we observe and enact the relational dynamics of power and control.